\chapter{Valutazione costi e ridondanza}
\paragraph{Attributo ridondante} consideriamo due operazioni sull'entità Workshop: 
\begin{itemize}
    \item \textbf{[OP-A]} Verifica posti disponibili (100 volte/mese)
    \item \textbf{[OP-B]} Iscrizione studenti al workshop (40 volte/mese) 
\end{itemize}
Introduciamo dunque \texttt{Numero\_partecipanti} al fine di velocizzare le operazioni di lettura all'interno del database, calcoliamo se mantenere l'attributo ridondante sia conveniente o meno: 
\paragraph{Con attributo ridondante} 
\textbf{[OP-A]}: 100 letture al mese
\newline
\textbf{[OP-B]}: 40 letture al mese + 40 scritture sull'attributo ridondante
\textbf{Totale}: $100 + 40 + (40 \times 2) = 260$ operazioni/mese 
\footnote{Considerando ogni scrittura come equivalente a due letture in costo computazionale}
\paragraph{Senza attributo ridondante} 
\textbf{[OP-A]} Ipotizzando una media di 8 partecipanti per ogni workshop avremo: 
\begin{enumerate}
    \item 1 lettura su Workshop
    \item 1 lettura su calendario\_workshop (nuova relazione necessaria)
    \item 8 conteggi per i partecipanti 
\end{enumerate}
Dunque $10 \times 100 = 1000\ operazioni/mese$ 
\newline
\newline
\textbf{{OP-B}} 40 scritture in Workshop $\longrightarrow 160\ operazioni/mese$
\newline
\textbf{Totale}: $1000 + (40 \times 2) = 1160$ operazioni/mese 
\paragraph{Conclusione} Conviene mantenere l'attributo ridondante considerato il discreto risparmio in termini di costo mensile per le operazioni